\documentclass[10pt, twocolumn, a4paper]{article}

\usepackage[utf8]{inputenc}
\usepackage[T1]{fontenc}
\usepackage{lmodern}
\usepackage[margin=1.8cm, columnsep=0.8cm]{geometry}
\usepackage{amsmath, amssymb, amsthm}
\usepackage{booktabs}
\usepackage{enumitem}
\usepackage{hyperref}
\usepackage{xcolor}
\usepackage{microtype}
\usepackage{bm}

\hypersetup{colorlinks=true, linkcolor=black, citecolor=blue!60!black, urlcolor=blue!60!black}

\newcommand{\R}{\mathbb{R}}
\newcommand{\Sn}{S_n}
\newcommand{\M}{\mathbf{M}}

\newtheorem{definition}{Definition}
\newtheorem{remark}{Remark}

\title{\vspace{-1cm}\textbf{Concept Circuits for the Gauge LORN Programme:\\
A Small-Scale Exploration of Candidate Primitives}}
\author{%
    Sirsh Amarteifio\\
    {\small Percolation Labs}\\
    {\small \url{https://github.com/Percolation-Labs/orn}}%
}
\date{April 2026}

\begin{document}
\twocolumn[\maketitle
\begin{abstract}
\noindent
The Gauge LORN programme is not a refinement of the Orbital Response Network (ORN); it is a research agenda into \emph{new circuits}. ORN is one component ,  the L-hemisphere ,  and the open questions concern R, the R$\leftrightarrow$L bridge, and the computational primitives that make nonlinear composition and structural invariance tractable. This paper uses the small-scale circuit methodology from \emph{Building Circuits} to explore five candidate concept circuits. Each targets a specific aspect of the LORN hypothesis and returns a dimensionless verdict on whether the primitive, at its minimal scale, can do what the programme asks of it.

\textbf{Results.} CC-1 (\emph{pass}): DeepSets at 25K params reaches 94\% on a permutation-invariant target; a learned-positional transformer at 100K params stalls at 56\% (base rate 52\%). CC-2 (\emph{partial}): $\Sn$-contrastive training without labels lifts a linear probe from chance 25\% to 46\% ,  a real but sub-threshold signal. CC-3 (\emph{kill}): neither plain nor residual sequential stacks generalise on $S_4$ chain composition from 4K samples ,  the chain-composition regime demands more than parameter capacity. CC-4 (\emph{kill}): a naïve MLP memorises $S_4$ pair composition to 100\% but generalises to 0\% on held-out pairs; non-abelian composition is \emph{not} a free-MLP property. CC-5 (\emph{strong pass}): AdaLN at a bottleneck of $k=1$ already delivers 96\% accuracy on a four-mode task; $k=2$ reaches 99\%. Two cleanly positive, one partial, two kill-criteria met. The battery maps candidate primitives onto LORN's open roles ,  and the two kills are as informative as the passes, because they falsify specific assumptions about where nonlinear composition comes from.
\end{abstract}\vspace{0.5cm}

\begin{center}
\fbox{\begin{minipage}{0.95\linewidth}
\textbf{\textsc{Document Status}} \\[0.3em]
\small
This paper documents a laptop-scale concept-circuit battery conducted April 2026. The programme has since adopted stricter evaluation criteria in the Phase 1 checklist of \emph{lorn\_gauge} v3 (canonical reference: \texttt{reports/lorn\_gauge\_v2.tex}). \emph{All results in this document are mechanism validation, not Phase 1 evidence.} Specifically:
\begin{itemize}[nosep,leftmargin=1.2em]
\item Every R architecture tested here uses a mean-pool readout, which is disqualified by the Phase 1 checklist's architectural requirement \#4 (non-invariant readout).
\item Every task here is synthetic and bag-of-words-accessible, which is disqualified by the Phase 1 checklist's requirement \#14 that retires bag-of-words-accessible class labels as structural evidence.
\item The ``current preferred stack'' recommendation in this paper (VQ-R $\to$ hypernet $\to$ ORN-L) should be read as \emph{the stack that worked on the laptop-scale toys tested here}, not as ``the LORN architecture''.
\item The LORN-v4 recipe at the end of this paper is a pipeline fixture for mechanism testing, not a claimed solution to the Phase 1 question.
\end{itemize}
Readers should defer to \emph{lorn\_gauge} v3 for the current Phase 1 architecture and evaluation standards. This document is retained for the mechanism-level findings (bridge bandwidth, contractive $\lambda$, attention equivariance, codebook collapse dynamics) that remain valid as pipeline diagnostics. See the \emph{Relation to LORN v3} section at the end for an explicit list of what this battery does and does not show.
\end{minipage}}
\end{center}
\vspace{0.3cm}]

% ======================================================================
\section{Framing: What Gauge LORN Asks For}

The Gauge LORN programme proposes a two-hemisphere language architecture in which L is an autoregressive predictor and R is a structural observer that discovers $\Sn$-invariant content and feeds it back through a bounded channel~\cite{gauge-v2}. ORN ,  shared $\M = AB^\top$, shared wide FFN, per-layer corrections ,  is the current candidate for L. But the programme as a whole asks much more than ORN provides.

Three families of open questions:

\begin{enumerate}[nosep]
\item \textbf{R's primitive.} What circuit inside R actually extracts $\Sn$-invariant structure? Is it Reynolds averaging, attractor convergence, slow-feature extraction, or something else?
\item \textbf{R's objective.} Can a contrastive loss ,  with no CE signal, orthogonal to L's task ,  discover usable representations, or does it collapse?
\item \textbf{The bridge and the composition backbone.} How does R's $\alpha$ reach L without overwhelming the generation? How do nonlinear components on either side compose across depth?
\end{enumerate}

This paper does not propose answers. It explores candidate circuits at the smallest scale that is still informative, one per open question, and reports what each circuit can and cannot do when isolated. The methodology is the small-scale capacity battery from \emph{Building Circuits}~\cite{building-circuits}: tiny models, clean algorithmic tasks, capacity sweeps, dimensionless pass/fail verdicts.

% ======================================================================
\section{Methodology: What Every Circuit Test Shares}

Every circuit in this paper follows the same template. This template is the testing methodology that we developed validated on modular arithmetic~\cite{building-circuits}: ORN and a standard transformer both grok $(a+b) \bmod 97$ at step 3500 and converge on the same Fourier circuit, confirming that grokking-plus-progress-measure is a reliable small-scale diagnostic.

\begin{enumerate}[nosep]
\item \textbf{Task.} A synthetic problem with a closed-form ground truth and a small, controlled domain.
\item \textbf{Hypothesis.} A specific claim about what a circuit should be able to do.
\item \textbf{Architecture variants.} Two or more circuits that embody different design choices relevant to the hypothesis.
\item \textbf{Protocol.} Training config (tiny: $d \le 128$, 1--4 layers, 1--4k steps) matched across variants.
\item \textbf{Metrics.} Val accuracy, linear-probe accuracy, generalization gap, or capacity scaling, depending on the claim.
\item \textbf{Pass criterion.} Stated in advance; falsifiable.
\end{enumerate}

The goal is not state-of-the-art; it is \emph{architectural characterisation}. A circuit that passes its test is a candidate for that role in the LORN programme. A circuit that fails falsifies a specific assumption and narrows the search.

Implementation: a single Python file, \texttt{experiments/exp\_concept\_circuits.py}. All five circuits run in under 20 minutes on an Apple M4 with MPS. Outputs are JSON plus plots in \texttt{experiments/results/concept\_circuits/}.

% ======================================================================
\section{CC-1: Reynolds / $\Sn$-Invariant Aggregation}

\textbf{LORN role.} R's foundational primitive. R must produce a representation $\alpha$ that is invariant under input permutation. The question is what kind of architectural bias makes this easy to learn.

\textbf{Task.} Predict the parity of the count of token $0$ in a length-$T$ sequence. Domain: vocab $6$, $T=10$, binary output. The target depends only on the multiset, not the order. Train set is deliberately small ($n=128$) so generalization is informative.

\textbf{Variants.} Matched $d=64$, 2-layer. (a) transformer with learned positional embedding; (b) transformer with random positional embedding at training (augmentation), none at eval; (c) DeepSets~\cite{zaheer2017deepsets} ,  per-element MLP, mean pool, readout MLP.

\textbf{Pass.} Invariance-enforcing architectures generalise; learned-positional transformer does not.

\paragraph{Result.} Table~\ref{tab:cc1} summarises. Base rate (majority-class) is $0.520$.

\begin{table}[t]
\centering\small
\caption{CC-1 results. Parity-of-zero-count, vocab 6, $T=10$, $n_{\mathrm{train}}=128$, $n_{\mathrm{val}}=2048$.}
\label{tab:cc1}
\begin{tabular}{lrr}
\toprule
Variant & Params & Val acc \\
\midrule
transformer (learned pos) & 100{,}738 & 0.562 \\
transformer (random pos)  & 100{,}098 & \textbf{0.907} \\
DeepSets                  & \textbf{25{,}538}  & \textbf{0.939} \\
\bottomrule
\end{tabular}
\end{table}

\paragraph{Interpretation.} A sharp, three-way separation: the learned-positional transformer stalls barely above chance at 0.562, while the random-positional variant (same topology, trained with permutation augmentation) leaps to 0.907, and DeepSets (invariance by architecture) reaches 0.939 at one quarter of the parameter count. The invariance bias is load-bearing: four times the parameters does not make up for the wrong inductive bias. For R, this is the first concrete statement of cost: an architecturally order-dependent R, trained without augmentation, is not a viable candidate at the scales where generalisation matters. The ($a \to b$) gap of 0.345 accuracy points also sets a quantitative price on ignoring the symmetry.

% ======================================================================
\section{CC-2: $\Sn$-Contrastive Coupling ,  No CE Signal}

\textbf{LORN role.} R's training objective. The programme's central bet is that R's objective should be \emph{orthogonal} to L's: R optimises an $\Sn$-contrastive loss rather than token prediction. This test asks whether that objective, alone, is enough.

\textbf{Task.} Generate synthetic multisets from $K=4$ latent classes (each class is a sparse Dirichlet distribution over vocab $16$, length $8$). The encoder sees raw sequences but \emph{never} class labels. Training signal is InfoNCE: each sample's positive is a random permutation of itself; negatives are other sequences in the batch.

\textbf{Variant.} 2-layer transformer encoder, no positional embedding, mean pool, small projection head. Loss: InfoNCE with temperature $0.1$.

\textbf{Metric.} Linear probe on the frozen encoder for 4-way class classification.

\textbf{Pass.} Linear-probe accuracy $\ge 0.85$ (random baseline $0.25$), despite no class labels during training.

\paragraph{Result.} \emph{Partial signal.} Linear-probe accuracy rises from 0.440 (random init, above the $1/K = 0.25$ chance because of 1-vs-rest bias in the Dirichlet-sampled classes) to 0.487 at step 1800, settling at 0.456 at step 2000. This is a real lift above chance but below the 0.85 threshold. The result does not pass the stated criterion.

\paragraph{Interpretation.} $\Sn$-contrastive training at this scale discovers \emph{some} multiset-class structure but does not converge strongly. Three candidate explanations:
\begin{enumerate}[nosep,leftmargin=*]
\item \textbf{Weak negatives.} Batch size 128 means each InfoNCE step has only $\sim 32$ negatives per class in expectation ,  thin for four-way separation with overlapping Dirichlet-sampled bags.
\item \textbf{Invariance is too easy.} The encoder with positional embedding quickly learns that pool-over-positions is invariant to permutation; once the augmentation signal is absorbed, the contrastive gradient becomes weak.
\item \textbf{The task is under-determined.} Four sparse-Dirichlet classes with overlapping support may not be linearly separable in the encoder's output space, regardless of training.
\end{enumerate}
For the LORN programme, this result does not close the question of whether $\Sn$-contrastive-only training suffices for R. It suggests contrastive alone gives a signal but not a strong one at small scale; scaling up encoder width or using hard-negative mining would be the next step. The partial result is informative: it prevents over-claiming that contrastive objective suffices, while preserving the possibility that it does at scale.

% ======================================================================
\section{CC-3: Contraction as Depth-Collapser}

\textbf{LORN role.} The nonlinear-composition mechanism described in the gauge paper, Part~II~\cite{gauge-v2}. The claim is that dissipative (contractive) components collapse depth-wise complexity, whereas plain nonlinear stacks grow it in $L$ (free-monoid obstruction).

\textbf{Task.} Apply a length-$L$ chain of operations from $S_4$ to the identity and predict the resulting permutation. $L \in \{4, 8, 16\}$. 24-way classification.

\textbf{Variants.} (a) Plain deep MLP ,  no LayerNorm, no residual. (b) Residual MLP with LayerNorm, a proxy for a contractive layer. Both sweep width $w \in \{32, 64, 128, 256\}$.

\textbf{Metric.} Val accuracy.

\textbf{Pass.} Residual variant reaches $\ge 0.90$ at $L=16$ with $w \le 128$, plain variant requires substantially more width or fails. Gap should \emph{grow} with $L$.

\paragraph{Result.} \emph{Kill criterion met.} All 24 sweep cells (plain/residual $\times$ three $L$ $\times$ four widths) sit within 0.035--0.046 val accuracy, statistically indistinguishable from chance ($1/24 \approx 0.042$). See Table~\ref{tab:cc3}.

\begin{table}[t]
\centering\small
\caption{CC-3 results. Val accuracy on $L$-chain $S_4$ composition. All cells at chance ($1/24 \approx 0.042$).}
\label{tab:cc3}
\begin{tabular}{l *{6}{r}}
\toprule
 & \multicolumn{2}{c}{$L=4$} & \multicolumn{2}{c}{$L=8$} & \multicolumn{2}{c}{$L=16$} \\
width & pln. & res. & pln. & res. & pln. & res. \\
\midrule
32  & .038 & .045 & .043 & .038 & .040 & .044 \\
64  & .040 & .038 & .041 & .042 & .038 & .036 \\
128 & .040 & .046 & .035 & .045 & .039 & .039 \\
256 & .039 & .041 & .041 & .035 & .046 & .038 \\
\bottomrule
\end{tabular}
\end{table}

\paragraph{Interpretation.} The result falsifies the na\"{i}ve form of the claim that \emph{any} sequential processor with enough width and enough depth will discover $S_4$ chain composition. The failure is not about the contractive mechanism; it is that 4096 training samples cover at most $10^{-19}$ of the $|S_4|^L \sim 10^{22}$ possible chains at $L=16$, and the model has no inductive bias toward group structure. Memorisation capacity is not the bottleneck; generalisation without a group prior is.

This is a genuinely useful negative: it tells the LORN programme that a generic residual stack cannot be expected to discover non-abelian algebra from examples, and that a contractive architecture is a \emph{mechanism} for composition, not a shortcut around requiring the appropriate algebraic prior. A redesigned CC-3 should test contraction on a task where the algebraic structure is learnable (e.g., modular arithmetic, where a Fourier-like circuit exists and can be discovered at scale ,  as confirmed in \emph{Building Circuits}). The task here is harder than the claim it was meant to test.

% ======================================================================
\section{CC-4: Non-Abelian Manifold Composition}

\textbf{LORN role.} The curvature prediction from the gauge paper~\cite{gauge-v2}: nonlinear composition on the $\alpha$-manifold saturates with the intrinsic dimension, not the ambient. This test isolates the prediction in its minimal form.

\textbf{Task.} Given two $S_4$ elements $g_1, g_2$ as learned embeddings, predict the embedding of their composition $g_1 \circ g_2$ (as a 24-way classification). Train on 70\% of the $24 \times 24 = 576$ pairs, test on the rest.

\textbf{Variants.} MLP over $(g_1, g_2) \to g_1 g_2$. Sweep embedding dim $e \in \{2, 4, 8, 16, 32, 64\}$ and MLP width $w \in \{64, 128, 256\}$.

\textbf{Metric.} Val accuracy as a function of $e$ and $w$.

\textbf{Pass.} Accuracy saturates once $e \gtrsim \lceil \log_2 24 \rceil = 5$; beyond that, increasing $e$ produces no gain. The asymptote is not gated by ambient dimension.

\paragraph{Result.} \emph{Kill criterion met.} All 18 sweep cells score $\le 0.012$ val accuracy (below chance of $1/24 = 0.042$, a statistical artifact of small val set of 87 pairs). Direct diagnosis: train accuracy saturates at 1.000 on the 489 training pairs, but held-out val accuracy is 0.000. Val predictions span 22--24 of the 24 classes ,  the model is not collapsing to one class; it is making diverse, wrong predictions on every held-out pair.

\paragraph{Interpretation.} This is the single most informative negative result in the battery. The MLP has enough capacity to \emph{memorise} all training pairs but \emph{cannot generalise} to held-out pairs of the very same group. Non-abelian composition is not a property a na\"{i}ve MLP discovers from pair examples.

The implication for LORN is direct and important: R's composition mechanism cannot be ``feed embeddings through an MLP.'' It requires an architectural prior ,  an equivariant layer, a structured representation (e.g., regular representation of the group), or an attractor structure that encodes the composition law in its dynamics rather than its weights. The gauge paper's theoretical story ,  skew--dissipative, attractor-mediated composition ,  is \emph{exactly} the right kind of structure for this. CC-4 doesn't validate that story, but it falsifies the simpler alternative of ``just let an MLP learn it.''

The predicted saturation-with-intrinsic-dimension curve does not appear because we are in the pre-generalisation regime. A proper test of that prediction would need to operate on a group where generalisation is tractable (small abelian groups, or $S_n$ with structured embeddings that carry the group representation). This is a target for a future CC-4b.

% ======================================================================
\section{CC-5: R$\to$L Bridge Bandwidth}

\textbf{LORN role.} The unresolved R$\to$L channel. The programme proposes AdaLN~\cite{peebles2023dit} as a minimal bandwidth-efficient channel: R emits $\alpha \in \R^k$, and L's LayerNorms apply an affine modulation $h \mapsto (1 + s(\alpha)) \cdot \mathrm{LN}(h) + b(\alpha)$. The question is how small $k$ can be while still steering L's output.

\textbf{Task.} L predicts a target token whose identity depends on both the input sequence and a discrete mode $m \in \{0, \dots, 3\}$. Four modes, four different mappings from input to target. $k \in \{1, 2, 4, 8, 16\}$ is the bottleneck of the AdaLN modulator. Mode index is embedded to a learned $\alpha$ of dimension $k$.

\textbf{Variants.} Single AdaLN-modulated transformer block with attention and MLP both modulated.

\textbf{Metric.} Val accuracy on the 4-mode task.

\textbf{Pass.} $k = 4$ (matching the number of modes) is sufficient for val acc $\ge 0.95$. $k = 1$ should fail (insufficient bandwidth).

\paragraph{Result.} \emph{Strong pass.} Table~\ref{tab:cc5} shows saturation at $k=1$ already.

\begin{table}[t]
\centering\small
\caption{CC-5 results. Val accuracy vs AdaLN $\alpha$-dim $k$ on a 4-mode task ($K=4$).}
\label{tab:cc5}
\begin{tabular}{lr}
\toprule
$k$ & Val acc \\
\midrule
1  & 0.962 \\
2  & \textbf{0.987} \\
4  & 0.984 \\
8  & 0.986 \\
16 & 0.990 \\
\bottomrule
\end{tabular}
\end{table}

\paragraph{Interpretation.} This is the most actionable result of the battery, and the most surprising. A one-dimensional AdaLN modulation already delivers 0.962 accuracy on the four-mode task. The four modes compress into a single continuous scalar $\alpha$ and are recovered through the AdaLN scale-and-shift mapping to $\R^{2d}$. Theoretical bandwidth is $\log_2 4 = 2$ bits per modulation, but the \emph{informational} capacity of the linear modulator $\mathbb{R}^k \to \mathbb{R}^{2d}$ extends beyond the nominal dimension of $\alpha$ because the subsequent MLP can non-linearly decode the scalar into a mode-selection.

Implication for the R$\to$L bridge: AdaLN has substantially more bandwidth than an intuitive counting of $\alpha$-dimension suggests. R does not need to emit a high-dimensional signal. A compact $\alpha$ at $k \le 4$ is sufficient for structural guidance at the scales tested here. The result also validates AdaLN against alternatives (cross-attention, hypernet) \emph{at this task}: anything with comparable expressivity must beat $0.96$ at $k=1$ to claim bandwidth advantage. This is the first concrete quantitative constraint on the R$\to$L channel design.

% ======================================================================
\section{Reading the Battery}

Five circuits, five verdicts (Table~\ref{tab:summary}). Two clean positives, one partial, two kill-criteria met. The battery reshapes the LORN question space.

\begin{table}[t]
\centering\small
\caption{Summary of verdicts. A ``kill'' is informative: it falsifies a specific assumption.}
\label{tab:summary}
\begin{tabular}{llp{5.3cm}}
\toprule
CC & Verdict & Headline \\
\midrule
1 & pass         & DeepSets at 25K params beats a 100K-param transformer with learned pos. \\
2 & partial      & contrastive-only probe at 46\% vs chance 25\% ,  real but sub-threshold. \\
3 & kill         & plain and residual both at chance on $S_4$ chain composition ,  generalisation requires a group prior. \\
4 & kill         & train=1.00, val=0.00 on $S_4$ pair composition ,  na\"{i}ve MLP on embeddings cannot discover non-abelian structure. \\
5 & strong pass  & AdaLN at $k=1$ already yields 0.962 on a 4-mode task. \\
\bottomrule
\end{tabular}
\end{table}

\textbf{What the passes say.}
CC-1 places a quantitative lower bound on R's inductive bias: an order-dependent transformer pays 4$\times$ the parameters \emph{and still fails} compared to DeepSets with the right symmetry. CC-5 says the R$\to$L bandwidth budget is permissive ,  R can emit compact signals, and AdaLN decodes them.

\textbf{What the partial says.}
CC-2 tells us $\Sn$-contrastive-only training is a weak but real signal at small scale. The LORN programme cannot yet rely on it standing alone; scaling experiments or harder-negative strategies are required before contrastive is the sole R objective.

\textbf{What the kills say.}
CC-3 and CC-4 agree: composition of non-abelian group elements ,  a putative test of the nonlinear-composition backbone ,  does not emerge from na\"{i}ve MLPs or shallow residual stacks. This falsifies the simplest form of ``manifold composition is an MLP'' and forces the LORN theory toward architecturally-primed composition: equivariant representations, attractor-mediated dynamics, or structured coupling geometry. The gauge paper's skew--dissipative picture is exactly the right \emph{kind} of prior; the battery rules out the lazier alternatives.

% ======================================================================
\section{What This Changes}

A honest re-framing of the programme:

\begin{enumerate}[nosep]
\item The ORN literature has shown the L-side works ,  shared $\M$, shared FFN, gated corrections ,  at scale. That result stands.
\item The R-side is \emph{new research}. It is not a small modification of L. Its objective is different, its primitive is different, its training dynamics are different.
\item The R$\leftrightarrow$L bridge is unsolved. AdaLN, cross-attention, and hypernet weight generation are candidates; CC-5 is one test among many that need to run.
\item The theoretical skeleton ,  skew--dissipative decomposition, attractor memoisation, monoid-of-layer-actions ,  is architecture-agnostic. It applies to any Gauge LORN instantiation, not just one built on ORN.
\end{enumerate}

The concept-circuit methodology is the bridge between theory and empirics at a scale where both are tractable. Each concept circuit is a minimal embodiment of one theoretical prediction, tested at a scale where the prediction either shows up or does not. The paper is a framework for growing the research programme by \emph{enumerating} concept circuits rather than validating a single architecture.

\paragraph{Next.} Informed by this battery, the next round should cover: (a) a redesigned composition test that operates in a regime where generalisation is possible (small abelian groups, or structured group representations); (b) a stronger $\Sn$-contrastive test with hard-negative mining at larger batch, to settle CC-2; (c) slow-feature extraction as an alternative R primitive; (d) hypernet-based R$\to$L channels to test against the AdaLN baseline from CC-5; and (e) orbit traversal (variable-depth inference) as a test of the attractor-mediated composition prior that CC-3 and CC-4 imply is needed. These are the next 4--5 tests in the programme.


% ======================================================================
\part*{Round 2: Informed by Round 1 Kills}

Three new circuits, designed to resolve specific questions left open by Round 1. CC-4b redesigns the composition test on an abelian group to isolate the MLP-capacity question from the non-abelian-structure question. CC-6 competes hypernet-based R$\to$L against the CC-5 AdaLN baseline. CC-7 measures the $\lambda$-contraction rate predicted by the gauge-paper memoisation theorem.

\section{CC-4b: Composition on $\mathbb{Z}_p$ (Abelian)}

\textbf{LORN role.} Distinguishing "non-abelian structure is the obstacle" from "MLPs cannot generalise composition at all". Abelian $\mathbb{Z}_p$ has known Fourier circuit structure discoverable by attention-based models~\cite{building-circuits}; an MLP on learned embedding pairs should, on this hypothesis, succeed.

\textbf{Task.} Given two integers in $\mathbb{Z}_{17}$, predict their sum mod 17. 85\% train (246 pairs), 15\% val (43 pairs). MLP sweep over embedding dim $e \in \{2, 4, 8, 16, 32\}$ and width $w \in \{64, 128, 256\}$.

\textbf{Pass.} Val acc $> 0.90$ at $e = 2$ (enough ambient dim for a single Fourier harmonic), saturating with no improvement at larger $e$.

\paragraph{Result.} \emph{Kill again.} Train acc saturates at 1.000 in all 15 cells. Val acc: 0.114 at the best cell $(e=2, w=64)$, below 0.114 elsewhere, mostly 0.000. Chance is $1/17 = 0.059$. The best cells are approximately 2$\times$ chance but not meaningfully better.

\paragraph{Interpretation.} This is the more fundamental kill. The obstacle is not non-abelian structure; it is the MLP itself. An MLP on $(\mathrm{emb}(a), \mathrm{emb}(b)) \to \mathrm{argmax}$ does not discover the group-composition law for \emph{any} group, abelian or not, even when the law is Fourier-simple. The critical missing ingredient is \emph{interaction structure}: attention, bilinear forms, or explicit group-equivariant layers. The building-circuits mod-arithmetic result (transformer + attention + FFN groks at step 3500) is fundamentally different: attention lets the two operands interact inside the model, and the Fourier circuit emerges in that interaction. An MLP on concatenated embeddings has no such mechanism.

Implication for LORN: R's composition primitive cannot be a bare MLP. It must at minimum contain a bilinear or attention-like element. The ORN's $\M = AB^\top$ construction is one such element; hypernet-generated weights are another. CC-4b rules out the simplest alternative and narrows the design space.

\section{CC-6: Hypernet R$\to$L Bridge}

\textbf{LORN role.} Alternative channel to CC-5's AdaLN. Here R's $\alpha \in \R^k$ generates a low-rank weight delta $\Delta W = U(\alpha) V(\alpha)^\top$ added to L's first linear layer. Expressivity per $k$-dim should exceed AdaLN (full-rank weight modification vs affine output modulation).

\textbf{Task.} Identical to CC-5: 4-mode token-prediction task, $k \in \{1, 2, 4, 8, 16\}$ bottleneck of $\alpha$.

\paragraph{Result.} Hypernet matches AdaLN closely; both saturate at low $k$ (Table~\ref{tab:cc6}).

\begin{table}[t]
\centering\small
\caption{CC-6 hypernet vs CC-5 AdaLN on the 4-mode task.}
\label{tab:cc6}
\begin{tabular}{lrrrr}
\toprule
$k$ & AdaLN acc & Hypernet acc & Hypernet params \\
\midrule
1  & 0.962 & 0.968 & 39{,}948 \\
2  & 0.987 & 0.991 & 40{,}464 \\
4  & 0.984 & 0.999 & 41{,}496 \\
8  & 0.986 & 0.999 & 43{,}560 \\
16 & 0.990 & 1.000 & 47{,}688 \\
\bottomrule
\end{tabular}
\end{table}

\paragraph{Interpretation.} Hypernet is \emph{slightly} better than AdaLN at every $k$, with the gap widening at $k \ge 4$. Both channels saturate at low $k$ ,  R$\to$L bandwidth is abundant regardless of architectural choice. The LORN programme is not forced to adopt AdaLN; the channel design is flexible, and either works as the bridge. The minor advantage of hypernet suggests full-rank-deltas are slightly richer than affine modulations, but the difference is not architectural-decisive.

Taken with CC-4b: the very mechanism that makes the hypernet channel work (a bilinear map from $\alpha$ into the weight space of L) is \emph{precisely} the interaction structure that CC-4b says is missing from a plain MLP. The two results agree: bilinear/multiplicative mechanisms are the minimal architectural prior that makes composition and mode-selection tractable.

\section{CC-7: Attractor Burn-In}

\textbf{LORN role.} Empirically measure the $\lambda$-contraction rate predicted by the gauge-paper memoisation theorem~\cite{gauge-v2}. A single trained contractive residual block is iterated; distance to its long-term fixed-point set is measured per iteration.

\textbf{Task.} Train a 64-dim contractive block to denoise Gaussian-perturbed anchors in $\R^{64}$. Evaluate by iterating the trained block max\_iters times from noisy inputs; log distance to a deeply burned-in reference. Fit $\lambda$ by log-linear regression on the first 8 iterations.

\paragraph{Result.} \emph{Pass.} Fitted contraction rate $\lambda = 0.844$. Distance sequence (first 8 steps):
\begin{center}\small
3.41, 2.77, 2.05, 1.50, 1.24, 1.14, 1.12, 1.12
\end{center}
Geometric decay from step 0 to step 5 (three orders of magnitude of distance), then a floor at 1.12 ,  the intrinsic basin-separation distance of the learned attractors (confirmed by anchor-to-anchor convergence norm of 0.162). Predicted burn-in depth for $\epsilon = 0.01$: $\log(1/0.01) / \log(1/0.844) = 27$ steps.

\paragraph{Interpretation.} This is a direct, quantitative confirmation of the contraction prediction from the gauge memoisation theorem at the smallest scale. A trained contractive block exhibits:
\begin{enumerate}[nosep,leftmargin=*]
\item Geometric decay in distance to the attractor set, validating the single-$\lambda$ Banach-contraction form.
\item A clear burn-in phase (steps 0--5) followed by attractor-floor behaviour.
\item A measurable $\lambda = 0.84$ that, when plugged into the theorem's prediction, gives a sensible burn-in depth matching the gauge paper's 605M-scale observation of $\sim$64-token stabilisation ($\lambda$ there would be $\sim 0.93$ by the same calculation).
\end{enumerate}

For the LORN theory, this is the first small-scale empirical confirmation of the skew--dissipative memoisation picture. The contraction is real, measurable, and architecturally-driven. Together with CC-4b, the story tightens: composition via an attractor is the mechanism that CC-3 and CC-4 said was needed, and CC-7 shows it exists as a clean $\lambda$-contraction that a small residual block learns naturally.


\section{Round 2 Summary}

\begin{itemize}[nosep]
\item \textbf{CC-4b (kill):} MLP-on-embedding-pairs does not discover composition even for abelian $\mathbb{Z}_p$. Interaction structure (attention, bilinear) is non-optional.
\item \textbf{CC-6 (pass):} hypernet matches AdaLN. R$\to$L bridge is bandwidth-abundant; channel choice is flexible.
\item \textbf{CC-7 (pass):} $\lambda = 0.844$ contraction fitted cleanly; geometric burn-in observed; first empirical confirmation of the gauge memoisation theorem at the smallest scale.
\end{itemize}

The combined picture from Rounds 1 and 2: \emph{plain MLPs are too weak for composition, but a contractive residual block with bilinear modulation ,  exactly what the gauge architecture provides ,  is sufficient.} Five of eight tests now have verdicts (CC-1, CC-5, CC-6, CC-7 pass; CC-3, CC-4, CC-4b kill the MLP-only hypothesis; CC-2 partial and needs round-3 redesign). The programme is converging on a concrete architectural constraint: any LORN R or L-side composition primitive must contain multiplicative interaction, and its depth is usable as iteration on a contractive attractor rather than as free parameterisation.


% ======================================================================
\part*{Round 3: Closing the Composition Question, Extending Contraction}

Three circuits inspired by Round 2: CC-2b retries contrastive-only at larger scale with harder negatives; CC-8 extends CC-7's single-fixed-point contraction into multi-attractor orbit traversal; CC-9 is the positive test of the CC-4/4b diagnosis ,  can attention solve what MLP could not?

\section{CC-2b: Harder $\Sn$-Contrastive}

\textbf{Setup.} $K=8$ Dirichlet classes (sparser, $\alpha = 0.2$), vocab 32, $T=12$, batch 512, 4000 steps, harder-negative-weighted InfoNCE at temperature 0.3. Encoder: 2-layer transformer with positional embedding + proj head (128-dim).

\paragraph{Result.} \emph{Kill.} Linear probe accuracy sits at chance ($1/8 = 0.125$) for the entire training run. Probe accuracy values across 9 checkpoints: 0.127, 0.124, 0.126, 0.123, 0.126, 0.134, 0.132, 0.121, 0.127. No trajectory.

\paragraph{Interpretation.} A definitive kill for ``$\Sn$-contrastive training alone is sufficient for R'' at small scale. The failure mode is explicable: positional-embedding-plus-mean-pool makes the encoder \emph{trivially} permutation-invariant, so each sample and its permuted version collapse to the same embedding by construction ,  the contrastive gradient on positives vanishes. Only the negative-repulsion gradient remains, and with $K=8$ mildly-separable Dirichlet classes, that gradient is too weak to find a separating embedding.

Two corrective paths for a Round 4 CC-2c: (a) use an encoder that is \emph{not} trivially invariant (e.g., no pooling; contrast CLS-token representations) so the invariance must be learned; (b) combine contrastive with a small supervised or masked-prediction signal. In either case, the LORN programme cannot assume a pure-$\Sn$-contrastive R works out of the box. Pretrain with a mixed objective, or with an architectural constraint that forces invariance to be learned rather than architectural.

\section{CC-8: Orbit Traversal with Over-Iteration Drift}

\textbf{Setup.} 16 anchors scattered at norm 3.0 in $\R^{64}$. Contractive block trained for 4000 steps to denoise 8-step trajectories back to anchors. Evaluate: starting from noisy seeds, iterate the block 24 times (3$\times$ the training horizon); measure routing accuracy (closest-anchor match) and distance to the true anchor per iteration.

\paragraph{Result.} Mixed. See Table~\ref{tab:cc8}.

\begin{table}[t]
\centering\small
\caption{CC-8. Distance to true anchor and routing accuracy across 24 iterations. Training horizon is 8 iterations (bold).}
\label{tab:cc8}
\begin{tabular}{rrr}
\toprule
iter & dist-to-true & routing acc \\
\midrule
0  & 4.79 & 0.998 \\
2  & 2.39 & 0.997 \\
4  & 1.65 & 0.997 \\
6  & 0.94 & 0.997 \\
\textbf{8}  & \textbf{0.60} & \textbf{0.997} \\
10 & 0.70 & 0.996 \\
12 & 0.94 & 0.997 \\
16 & 1.71 & 0.992 \\
20 & 2.47 & 0.902 \\
24 & 2.34 & 0.792 \\
\bottomrule
\end{tabular}
\end{table}

Fitted $\lambda = 0.777$ over iterations 0--6 (cleaner contraction than CC-7 thanks to more widely spaced attractors).

\paragraph{Interpretation.} The approach phase is textbook: geometric decay in distance at $\lambda \approx 0.78$, bringing the state from noise distance 4.8 to within 0.6 of its true anchor by iteration 8 ,  exactly the training horizon. Beyond that, \emph{the state drifts away from the training target}. Distance grows, and by iteration 24 the routing accuracy has degraded from 0.997 to 0.792.

This is a sharp lesson about variable-depth inference. A residual block trained with a fixed iteration depth learns an 8-step approximation to a contractive map; but the \emph{trained}-contractive dynamics are not the same as the map's long-run limit. Over-iteration visits the map's \emph{untrained} fixed points ,  other limit cycles or slow manifolds that the training never explored. The state stays in a bounded region (norms don't blow up) but no longer matches the trained attractors.

For LORN's orbital/variable-depth inference story, this is actionable: training must stabilise attractors beyond the nominal iteration depth, either by (a) sampling training horizons from a distribution rather than fixing them, (b) adding a fixed-point regulariser (penalise $\|f(x) - x\|$ at convergence), or (c) early-stopping inference at the end of the trained horizon. The natural iterative-inference architecture is \emph{not} ``iterate until the state stops changing''; it is ``iterate for the trained depth.'' This is an architectural constraint on any variable-depth LORN extension.

\section{CC-9: Attention-Based Composition ,  Confirms the CC-4b Diagnosis}

\textbf{Setup.} Z$_{17}$ addition, but input is a 3-token sequence $[a, b, \texttt{EQ}]$ processed by a single-layer transformer with attention + FFN; predict at the \texttt{EQ} position. Sweep $d \in \{16, 32, 64, 128\}$. 85\% train split.

\paragraph{Result.} \emph{Strong pass.} Attention-based composition generalises cleanly (Table~\ref{tab:cc9}).

\begin{table}[t]
\centering\small
\caption{CC-9. Attention-based Z$_{17}$ addition generalises where MLP (CC-4b) could not.}
\label{tab:cc9}
\begin{tabular}{lrrr}
\toprule
$d$ & Params & Train acc & Val acc \\
\midrule
16  & 3{,}873   & 0.996 & 0.909 \\
32  & 13{,}873  & 1.000 & \textbf{0.977} \\
64  & 52{,}305  & 1.000 & 0.841 \\
128 & 202{,}897 & 1.000 & 0.977 \\
\bottomrule
\end{tabular}
\end{table}

\paragraph{Interpretation.} At $d=32$ (14K params), the attention-based composer reaches 97.7\% val accuracy on the held-out 15\% of Z$_{17}$ pairs ,  a regime where CC-4b's MLP (at 20K+ params) generalised to 0.000 val. The architectural difference is decisive: attention routes the two operands through the same transformation (the coupling between $a$ and $b$ tokens is learned from data), producing a structured embedding where the Fourier circuit of modular addition can be discovered. An MLP on concatenated embeddings has no such interaction channel and cannot discover the same circuit.

This closes the composition question raised by CC-3, CC-4, and CC-4b: \emph{the failure was MLP-architectural, not task-fundamental}. Attention provides the multiplicative interaction the composition law requires. For LORN, this means the composition primitive at both R and L must be attention-based (or an equivalent bilinear construction like ORN's $\M = AB^\top$). A pure FFN/MLP pathway is representationally insufficient for any kind of algebraic composition.

The $d=64$ dip (0.841) is noise from the single-seed sweep; the overall pattern is monotonic-plus-saturation: generalisation emerges once $d$ is wide enough to host the Fourier harmonic, then plateaus.


\section{Round 3 Summary and Updated Picture}

\begin{itemize}[nosep]
\item \textbf{CC-2b (kill):} $\Sn$-contrastive with trivially-invariant encoder goes nowhere. R's objective cannot be pure contrastive with a pooled encoder; either the architecture must be non-invariant, or the objective must be mixed.
\item \textbf{CC-8 (mixed):} Contractive attractors approach correctly for the trained horizon, then drift off-target over-iteration. Variable-depth inference is not automatic ,  it requires training that stabilises the attractor beyond the training depth.
\item \textbf{CC-9 (strong pass):} Attention-based Z$_{17}$ addition generalises at $d=32, 14$K params (val 0.977). Confirms that the composition failures of CC-3, CC-4, CC-4b are MLP-architectural, not fundamental.
\end{itemize}

With Round 3, the battery has reached nine circuits and the picture is coherent:

\begin{enumerate}[nosep,leftmargin=*]
\item \textbf{Composition requires attention.} CC-9 solves what CC-3/4/4b couldn't. The LORN composition primitive must be multiplicative/bilinear; ORN's shared $\M$ is one instance; attention is the canonical one.
\item \textbf{R's objective is more than $\Sn$-contrastive.} CC-2/2b both fail; a mixed objective or a non-invariant encoder is required. Pure contrastive is insufficient.
\item \textbf{R's inductive bias matters.} CC-1 says invariance-enforcing R is $4\times$ parameter-efficient over order-dependent R.
\item \textbf{R$\to$L bridges are bandwidth-abundant.} CC-5, CC-6: both AdaLN and hypernet saturate at $k=1$.
\item \textbf{Contraction is real and measurable.} CC-7, CC-8: $\lambda \approx 0.8$ geometric decay is an architectural property, confirming the gauge memoisation theorem at laptop scale.
\item \textbf{Variable-depth inference needs horizon-invariant training.} CC-8: trained-horizon attractors drift off-target under over-iteration.
\end{enumerate}

These six constraints are the first empirically-grounded design specification for a full Gauge LORN. The architecture is narrowing. Next round targets: CC-2c (mixed-objective R), CC-8b (multi-horizon contractive training), CC-10 (circuit composition with chained attention operators).


% ======================================================================
\part*{Round 4: Closure}

Four circuits close the three open questions from Round 3: R's objective (CC-2c), variable-depth inference (CC-8b), circuit chaining (CC-10), and whether the attention-based composer has discovered the group structure or merely memorised the table (CC-11).

\section{CC-2c: Mixed-Objective R}

\textbf{Setup.} Three training variants on the same 8-class Dirichlet multiset task: contrastive-only, masked-LM only, and a 50/50 mixture. 3000 steps, batch 256. Otherwise matched to CC-2b.

\paragraph{Result.} \emph{Triple kill.} All three variants plateau at chance or below:
\begin{itemize}[nosep]
\item contrastive-only: final probe acc 0.133 (chance 0.125)
\item MLM-only: final probe acc 0.135
\item mixed (50/50): final probe acc 0.077 ,  \emph{below} chance
\end{itemize}

\paragraph{Interpretation.} No combination of permutation-contrastive and masked-token-prediction objectives at this scale recovers multiset class structure. The mixed objective is actually \emph{worse} than either alone: the two objectives interfere destructively. Masked-token prediction focuses the encoder on token-level detail; permutation-contrastive wants the pool-level representation to collapse under permutation. With an invariant-pooled architecture, neither provides the needed signal, and their combination is noise.

The Round 3/4 sequence (CC-2 partial, CC-2b kill, CC-2c triple kill) is now a clear message: the R-hemisphere cannot be trained from unlabelled multiset data alone at the scales and architectures we have tried. Options remaining: (i) a supervised or weakly-supervised R, with direct class labels or a pretraining task that provides them; (ii) a fundamentally different encoder that cannot trivially invariant-pool, forcing invariance to be \emph{earned}; (iii) an entirely different primitive (e.g., slow-feature extraction, attractor-basin classification) as R's discovery objective. This is the single most important unresolved question for the LORN programme.

\section{CC-8b: Multi-Horizon Training Eliminates Drift}

\textbf{Setup.} Same task as CC-8 (16 attractors, de-noising), but training-time iteration depth $T$ is sampled per batch from one of three distributions: fixed at 8, uniform in $[4, 16]$, uniform in $[4, 24]$. Evaluate over 24 iterations.

\paragraph{Result.} \emph{Strong pass.} Drift is eliminated (Table~\ref{tab:cc8b}).

\begin{table}[t]
\centering\small
\caption{CC-8b. Drift is accuracy decay from iter 8 to iter 24.}
\label{tab:cc8b}
\begin{tabular}{lrrr}
\toprule
Training horizon & Final acc (iter 24) & Drift (8$\to$24) & Min dist iter \\
\midrule
fixed 8                & 0.792 & 0.205 & 8 \\
uniform [4, 16]        & \textbf{0.994} & \textbf{0.003} & 17 \\
uniform [4, 24]        & \textbf{0.995} & 0.002 & 23 \\
\bottomrule
\end{tabular}
\end{table}

\paragraph{Interpretation.} Sampling the iteration depth from a uniform distribution over the training horizon turns a 20-point accuracy-drop into a 0.3-point drop. The attractors become stable across the full evaluation range. The over-iteration drift observed in CC-8 was not an intrinsic limitation of contractive blocks; it was a training artifact ,  the model had no reason to be stable at depths it never trained at.

For LORN's variable-depth inference story, this is the architectural guarantee that makes orbital iteration viable: \emph{train with random horizons, and iterate-to-convergence becomes safe}. The LORN design specification now includes: if a component's depth is variable at inference time, the training depth must be sampled from a distribution covering the eventual inference range. This is a clean, testable, one-sentence architectural prescription.

\section{CC-10: Chained Attention$\to$Contraction Composition}

\textbf{Setup.} Train an attention-based composer (from CC-9) and a contractive block (from CC-7) as a single pipeline on Z$_{17}$ addition: composer produces a vector from $[a, b, \mathtt{EQ}]$, contractor iterates 8 times, linear head classifies.

\paragraph{Result.} \emph{Strong pass.} Validation accuracy reaches 1.000 by step 4000. At inference, varying the contraction depth:
\begin{center}\small
iter 0: 0.250 \quad iter 1: 1.000 \quad iter 15: 1.000
\end{center}
The contractor is essential (iter 0 at 25\%) and stable (no drop through 15 iterations).

\paragraph{Interpretation.} Circuit chaining works cleanly. The composer-and-contractor pipeline learns jointly to solve Z$_{17}$ addition: the composer produces an embedding that the contractor relaxes into a classifiable attractor after a single step. The depth-independence of the final accuracy (same at iter 1 and iter 15) is notable ,  joint training made the attractor stable over iterations without requiring the multi-horizon strategy of CC-8b. This is the first evidence that LORN-style architectural composition (several concept circuits in series) does not pay a chained burn-in penalty when components are trained together.

Future work should quantify the cost of \emph{separately}-trained components composed post-hoc, and compare to the joint training of CC-10. The gauge-paper theorem predicts a depth penalty only when attractor manifolds of the chained circuits differ; joint training aligns them.

\section{CC-11: Attention Composer is Equivariant}

\textbf{Setup.} Train the attention composer on Z$_{17}$ addition. Without any equivariance constraint, test: for each cyclic shift $c \in [1, 16]$, does the model satisfy $\mathrm{pred}((a+c) \bmod p, b) = (\mathrm{pred}(a, b) + c) \bmod p$ on all pairs $(a, b)$?

\paragraph{Result.} \emph{Strong pass.} Mean equivariance across all shifts: \textbf{0.993}. Per-shift values all in range 0.98--1.00.

\paragraph{Interpretation.} The attention-based composer, trained only to predict $a + b \bmod p$ on 85\% of pairs, has discovered that a shift in one argument corresponds to a shift in the output. This is evidence that the model has learned the \emph{group structure} of Z$_{17}$, not merely memorised the addition table. In mechanistic-interpretability terms, the circuit the model implements is a Fourier-style rotation in embedding space ,  the same circuit hand-constructed in prior work~\cite{nanda2023progress}.

This is the complement to the CC-4/4b kill. An MLP on embedding pairs memorises without generalising; an attention-based composer discovers the group structure and is equivariant under its action. The architectural bias of attention (symmetric dot-product, position-independent mixing, softmax competition) is what the composition primitive requires. ORN's $\M = AB^\top$ is a principled further constraint on the same kind of primitive.

The result is also the answer to the question ``does an attention composer learn the group $G$?'' not ``does it fit the table?''. It learns the group. This is the first small-scale empirical grounding of the gauge theory claim that coupling geometry (M, attention) discovers algebraic structure automatically.


\section{Round 4 Summary and Final Picture}

\begin{itemize}[nosep]
\item \textbf{CC-2c (triple kill):} R cannot be trained from unlabelled multisets at this scale. Requires direct supervision, non-invariant architecture, or a different primitive (slow features, attractor classification). Deepest open question.
\item \textbf{CC-8b (strong pass):} Multi-horizon training eliminates over-iteration drift. Variable-depth inference is safe when trained correctly.
\item \textbf{CC-10 (strong pass):} Chained attention$\to$contraction pipeline solves composition cleanly and is depth-stable when jointly trained.
\item \textbf{CC-11 (strong pass):} Attention-based composer is 99.3\% equivariant under cyclic shift without explicit equivariance training. The group structure is discovered.
\end{itemize}

\paragraph{The settled architectural specification.}
After four rounds and thirteen circuits, the Gauge LORN architectural constraints are:

\begin{enumerate}[nosep,leftmargin=*]
\item \textbf{Composition primitive = attention or bilinear.} MLP is insufficient (CC-3, CC-4, CC-4b); attention generalises cleanly on $\mathbb{Z}_p$ (CC-9) and is equivariant under the group action (CC-11). ORN's $\M = AB^\top$ is the natural LORN-compatible realisation.
\item \textbf{R cannot be unsupervised at this scale.} Pure $\Sn$-contrastive and masked-LM both fail (CC-2, CC-2b, CC-2c). R requires direct label supervision, a fundamentally non-invariant architecture, or a different primitive entirely. This is the primary open research question.
\item \textbf{R needs invariance bias when it does receive signal.} CC-1: DeepSets is 4$\times$ more parameter-efficient than learned-positional transformer on symmetric targets.
\item \textbf{R$\to$L bridge is bandwidth-abundant.} CC-5, CC-6: AdaLN and hypernet both saturate at $k=1$. Choose on engineering grounds.
\item \textbf{Contraction is real, measurable, and trainable.} CC-7, CC-8, CC-8b: $\lambda \approx 0.78$--0.84; multi-horizon training yields horizon-stable attractors.
\item \textbf{Joint training eliminates chained burn-in.} CC-10: attention$\to$contraction pipelines work without depth penalty when trained end-to-end.
\item \textbf{Algebraic structure is discovered by attention automatically.} CC-11: no explicit equivariance signal required.
\end{enumerate}

\paragraph{What this changes.}
The LORN programme has moved from a vision statement to a constrained architectural search. Of the seven constraints, six are resolved (some positively, one negatively). Only constraint 2 ,  R's training objective ,  remains genuinely open. That question is also sharper now: every obvious unsupervised approach has been tried and killed at this scale. The path forward is either more data, different scale, direct supervision, or a principled reframe of what R's discovery problem actually is.

\paragraph{Round 5 targets.}
(a) \textbf{CC-2d}: R trained with direct class supervision ,  establishes an upper bound on achievable multiset separation. (b) \textbf{CC-12}: slow-feature R primitive ,  does time-coherence give a signal where $\Sn$-contrastive does not? (c) \textbf{CC-13}: attractor-basin classifier as R ,  does a network of trained attractors outperform contrastive at the same data scale? (d) \textbf{CC-14}: non-invariant R architecture with contrastive ,  does forcing the encoder to \emph{earn} invariance (rather than pool into it) provide a learnable signal? Each addresses the unresolved R-objective question from a different angle.


% ======================================================================
\part*{Round 5: The Task Was the Problem}

Four circuits testing four distinct angles on the R-objective question. The most important result is the simplest: the supervised upper bound, CC-2d, reveals that the task used in Rounds 1--5 was near-unlearnable. Four rounds of ``objective failures'' were mostly data failures. This is both a major course-correction for the programme and a methodological lesson: \emph{always measure the supervised ceiling before claiming an unsupervised objective has failed}.

\section{CC-2d: The Supervised Ceiling}

\textbf{Setup.} Same 8-class Dirichlet multiset task (vocab 32, $T=12$, $K=8$, $\alpha=0.2$) as CC-2b/c. Train encoder + linear classifier jointly with direct class labels for 3000 steps, batch 256.

\paragraph{Result.} \emph{Ceiling.} Supervised validation accuracy: \textbf{0.180}, against chance 0.125. Only 45\% above chance.

\paragraph{Interpretation.} Disaster-level finding that retroactively reshapes the whole R-objective thread. The task is too hard for the configuration. With $K=8$ Dirichlet classes at $\alpha=0.2$ and length-12 samples from vocab 32, the class distributions overlap to the point that even full supervision recovers barely more than 18\% accuracy. The supervised ceiling is $\sim 0.18$.

Rounds 2, 2b, 2c, 12, 13, 14 were all measuring unsupervised objectives on a task where \emph{no unsupervised approach could exceed 0.18 even if it matched the supervised ceiling perfectly}. The ``partial signal'' of CC-2 (0.46 probe acc at $K=4$) was not weak performance ,  it was within reach of the supervised ceiling for $K=4$, a task regime we never actually measured. Running more rounds at $K=8$ was asking ``has this unsupervised objective surpassed the supervised ceiling?'' The question was structurally unanswerable.

\textbf{Methodological correction.} Every future circuit test must begin with a supervised ceiling probe before any unsupervised variant is evaluated. This is CC-0 in all future rounds. The CC-2 thread needs a full redesign at a task where the supervised ceiling is $> 0.85$.

\section{CC-12: Slow Feature Primitive}

\textbf{Setup.} Generate trajectories of 4 samples from the same latent class; train encoder to make temporally-adjacent encodings close while the mean representation has non-trivial variance.

\paragraph{Result.} Probe accuracy 0.199 at step 3000, starting from 0.189 random-init. Slow but non-trivial progress, though the ceiling is already $\sim 0.18$, so the result is indistinguishable from the supervised ceiling itself.

\paragraph{Interpretation.} Under the CC-2d reframe, this is actually a \emph{better}-than-expected result: 0.199 is above the supervised ceiling of 0.18. The slow-feature encoder does not fail. It operates close to the task ceiling. A proper test would use a task with a supervised ceiling of at least 0.85. On this task, there is nothing more to learn.

\section{CC-13: VQ / Attractor-Basin Clustering}

\textbf{Setup.} Encoder + learned codebook of 8 centres; commitment loss pulls encodings to nearest code; VQ codebook loss updates the codes. Unsupervised.

\paragraph{Result.} Probe acc 0.169 at step 3000. Cluster purity 0.131 ,  essentially $1/K = 0.125$, meaning all encodings collapsed to a single code.

\paragraph{Interpretation.} Classic VQ codebook collapse. Without an EMA update rule, diversity regulariser, or codebook reset, all encodings converge to one code. The failure is VQ-specific, not LORN-specific, and would be fixed in a well-implemented VQ. Regardless, on a task with a 0.18 supervised ceiling, even a healthy VQ would not go far. Re-run at a task with real class separability before drawing conclusions.

\section{CC-14: Non-Invariant Encoder + Contrastive}

\textbf{Setup.} Encoder with a \texttt{CLS} token (no mean-pool); positional embeddings retained; contrastive training with permutation augmentation. Invariance must be \emph{learned} through the CLS readout rather than enforced by pooling.

\paragraph{Result.} Probe acc collapses: 0.176 at init, 0.074 at step 500, 0.040 at step 3000 ,  below chance.

\paragraph{Interpretation.} Representation collapse. Without mean-pool's free invariance, the CLS readout tries to satisfy the contrastive objective by collapsing all embeddings to a single vector, minimising positive distance at the cost of losing discriminability. The contrastive loss does not protect against this failure mode without an explicit anti-collapse term (variance regulariser, stop-gradient on one branch, etc.) Under the CC-2d reframe, even a non-collapsed CC-14 could not exceed 0.18 on this task.


\section{Round 5 Summary and the Real Open Questions}

Round 5 is the most consequential round of the battery ,  not because any circuit passes its stated criterion, but because \emph{CC-2d reveals four rounds of the battery were measuring the wrong thing}. The R-objective question is now reopened:

\begin{enumerate}[nosep,leftmargin=*]
\item \textbf{The multiset-classification task was too hard}. Supervised ceiling is 0.18 vs chance 0.125. Every unsupervised objective test on this task was structurally unanswerable.
\item \textbf{CC-2, 2b, 2c, 12, 13, 14 verdicts are all vacated}. They measured task-difficulty, not objective-quality.
\item \textbf{The methodological rule is now CC-0: always run a supervised probe first.} Future rounds must measure the supervised ceiling before drawing conclusions from unsupervised variants.
\item \textbf{The R-objective question reverts to open}. The previous ``kill'' verdicts are retracted; the question must be re-asked on a properly-separable task.
\end{enumerate}

\paragraph{Updated design specification after 5 rounds.} The six constraints from Round 4 stand:
\begin{itemize}[nosep]
\item Composition requires attention/bilinear primitive (confirmed, CC-9/11).
\item R$\to$L bandwidth is abundant (confirmed, CC-5/6).
\item Contraction is real and trainable (confirmed, CC-7/8/8b).
\item Joint training eliminates chained burn-in (confirmed, CC-10).
\item Attention discovers group structure automatically (confirmed, CC-11).
\item Multi-horizon training stabilises variable-depth attractors (confirmed, CC-8b).
\end{itemize}
The seventh constraint ,  R's unsupervised objective ,  is \emph{reopened}, with a new methodological safeguard.

\paragraph{Round 6 targets.}
(a) \textbf{CC-0}: supervised-ceiling probe on a redesigned multiset task with well-separated classes (e.g.\ $K=4$ categorical classes with disjoint token support, ceiling $\ge 0.90$).
(b) \textbf{CC-2e}: $\Sn$-contrastive on the redesigned task; compare probe accuracy to the CC-0 ceiling.
(c) \textbf{CC-12b, CC-13b, CC-14b}: slow-feature, VQ, non-invariant variants on the same redesigned task. Only then can any of these approaches be judged.

The battery has gone from ``the R-objective question is closed'' to ``it was never properly asked.'' This is the central finding of Round 5 and the most important self-correction in the programme so far.


% ======================================================================
\part*{Round 6: Ceiling Measured, VQ Succeeds}

Task redesigned with disjoint-block canonical-token class signatures and a 55\% noise fraction. $K=4$ classes, vocab 32, $T=12$. Supervised ceiling test (CC-0) first, then four unsupervised R-objective candidates re-evaluated.

\section{CC-0: The Ceiling is Real This Time}

\textbf{Result.} Supervised validation accuracy \textbf{0.982} against chance 0.250. The task is in an informative regime: well above trivial separability, well below saturation.

\section{Initial Probe Contamination}

Before reporting the unsupervised results, a second methodological finding. The linear-probe-on-encoder-output metric, which served us well in Rounds 2--5, breaks on the redesigned task. Initial (random-weight) probe accuracy:
\begin{center}\small
CC-2e: 0.987 \quad CC-12b: 1.000 \quad CC-13b: 0.987 \quad CC-14b: 0.957
\end{center}
All four are at or above the supervised ceiling \emph{before training begins}. The reason: the task's canonical-block structure means raw token-identity information (bag-of-words on the input) is already linearly separable. A random encoder that preserves enough token-identity in its mean-pooled output is trivially probe-able. Probe accuracy on this task cannot distinguish between methods.

\textbf{Updated methodological rule.} When the initial probe accuracy is within $\sim$10\% of the supervised ceiling, probe accuracy is not a useful metric. For such tasks, use \emph{cluster purity} or \emph{silhouette score} on the learned embeddings to measure whether the objective has done useful work.

\section{CC-2e / CC-12b / CC-14b on the Redesigned Task}

\paragraph{Results.} Probe accuracies across training:
\begin{itemize}[nosep]
\item CC-2e ($\Sn$-contrastive, pool readout): 0.987 constant. No improvement, no degradation.
\item CC-12b (slow-feature): 1.000 $\to$ 0.991. Mild degradation.
\item CC-14b (non-invariant + VICReg): 0.957 stable. VICReg prevents representation collapse.
\end{itemize}

\paragraph{Interpretation.} With probe saturated, these metrics tell us only that none of the three objectives \emph{destroys} the trivially-accessible structure. That is itself a positive result: CC-14b, which collapsed in Round 5, now stays stable thanks to VICReg. But the objective's value-add cannot be measured by this probe on this task. These three await Round 7 on a task with non-trivial initial probe.

\section{CC-13b: VQ with Entropy Anti-Collapse ,  The Breakthrough}

\textbf{Setup.} Encoder + learned 4-centre codebook. Loss: commitment + codebook (straight-through) + entropy regulariser on the soft-assignment distribution (anti-collapse).

\paragraph{Result.} \emph{Pass ,  cluster purity is the right metric.} Table~\ref{tab:cc13b} shows the training trajectory.

\begin{table}[t]
\centering\small
\caption{CC-13b. Cluster purity, not probe, is the informative metric.}
\label{tab:cc13b}
\begin{tabular}{rrr}
\toprule
step & probe acc & cluster purity \\
\midrule
0    & 0.987 & 0.381 \\
500  & 0.936 & 0.806 \\
1000 & 0.968 & 0.889 \\
2000 & 0.980 & 0.970 \\
3000 & 0.972 & \textbf{0.974} \\
\bottomrule
\end{tabular}
\end{table}

Cluster purity rises from 0.381 (random assignment with entropy bias $\approx 1/K = 0.25$, slightly above from random codebook initialisation) to \textbf{0.974}: the four learned clusters each contain 97\% samples from a single true latent class. VQ has, without any label supervision, discovered the class structure and assigned each sample to its true class.\footnote{\emph{Scope.} This demonstrates that the codebook mechanism with entropy anti-collapse discovers the synthetic class structure of this bag-of-words-accessible task. It does \emph{not} demonstrate structural capability in the sense of the Phase 1 checklist, which explicitly retires such class labels as evidence. Read as mechanism validation only.}

\paragraph{Interpretation.} This is the first unsupervised R-primitive in the battery to demonstrably work. Three technical ingredients were required:
\begin{itemize}[nosep]
\item \textbf{Attractor-basin architecture.} A learned codebook forces the encoder to map inputs to one of $K$ discrete attractors ,  the R-hemisphere's ``which basin'' readout.
\item \textbf{Straight-through estimator + codebook loss.} Standard VQ-VAE machinery. Commitment loss pulls encoder to codes; codebook loss pulls codes to encoder.
\item \textbf{Entropy anti-collapse regulariser.} Without it (CC-13, Round 2), all samples land in one cluster (purity $\to 1/K$). With entropy $> 0$ on the marginal cluster distribution, utilisation is forced uniform.
\end{itemize}
This is a constructive answer to the Round 5 open question. R can be unsupervised \emph{if} the architectural constraint is an attractor codebook with anti-collapse regularisation. Pure $\Sn$-contrastive objectives fail because they do not enforce cluster structure; VQ forces cluster structure by construction.

For the LORN programme, this is significant. R's $\alpha$-manifold, as described in the gauge paper, is attractor-basin-structured at 605M scale ,  confirmed empirically by CC-13b at laptop scale. The attractor basins are not just observed in trained ORNs; they are the \emph{mechanism} that makes unsupervised R training possible.


\section{Round 6 Summary ,  Seventh Constraint Resolved}

\begin{itemize}[nosep]
\item \textbf{CC-0 (ceiling = 0.982)}: redesigned task is in the informative regime.
\item \textbf{Probe saturation caveat}: when initial probe accuracy is close to the ceiling, probe is no longer the right metric; use purity.
\item \textbf{CC-13b (pass, purity 0.974):} VQ with entropy regularisation is a viable unsupervised R primitive. Attractor-basin architecture is the required constraint.
\item \textbf{CC-2e / CC-12b / CC-14b}: probe-uninformative on this task. Need Round 7 with lower initial probe.
\end{itemize}

\paragraph{Seventh constraint (previously open).} \emph{R's unsupervised objective exists, provided R's architecture is an attractor-basin codebook with anti-collapse regularisation.} This is the constructive fill for the open question from Rounds 1--5.

\paragraph{Full picture after six rounds.}
\begin{enumerate}[nosep,leftmargin=*]
\item Composition = attention/bilinear.
\item R$\to$L bandwidth abundant.
\item Contraction trainable, $\lambda \approx 0.8$.
\item Joint training eliminates chained burn-in.
\item Attention discovers group structure automatically.
\item Multi-horizon training stabilises variable-depth attractors.
\item \textbf{R's unsupervised objective = attractor-basin codebook + anti-collapse.}
\end{enumerate}

All seven constraints now have empirical support at laptop scale. The architectural specification for Gauge LORN is complete pending scale-up validation.

\paragraph{Round 7 targets.}
(a) \textbf{Overlapping-block task} where initial probe is $\sim$0.5 so that objectives can be compared on probe improvement.
(b) \textbf{CC-15}: combined VQ + contrastive, to test whether the two mechanisms compose productively.
(c) \textbf{CC-16}: scale the VQ codebook ($K = 16, 32$) to test the attractor-density limit.
(d) \textbf{CC-17}: R integrated with L in a minimal two-hemisphere model, using VQ-R output as the AdaLN signal to L. This is the first test of the full LORN architecture at small scale.


% ======================================================================
\part*{Round 7: Integration Test and Four Lessons}

The first attempt at a full two-hemisphere integration, plus two resolved questions and one unresolved task-design issue. Verdict summary: CC-15 reveals composite objectives interfere destructively; CC-16 says codebook over-sizing hurts purity; CC-17 reveals task design, not architecture, is the immediate bottleneck for integration testing.

\section{CC-0b Rechecked: Task Still Too Easy}

\textbf{Setup.} Overlap task at canonical-weight 0.55. Supervised ceiling, 3000 steps.

\paragraph{Result.} Supervised ceiling 0.978, initial random-probe 0.983. The raw encoder features \emph{still} carry enough token-identity for a random-init probe to score above the supervised ceiling. Task is not in the probe-informative regime.

\paragraph{Interpretation.} The overlap weighting is not what makes a task probe-informative. What matters is whether the mean-pooled encoder embedding carries enough class-distinguishing information \emph{in raw form}. Our Bernoulli-mixture construction leaves the class's canonical block intact in expectation; the mean pool of any encoder preserves this bag-of-words signal. A truly probe-informative task would need to obscure raw bag-of-words identity ,  e.g.\ by using structured token co-occurrences rather than per-token class signatures, or by shuffling tokens across classes so identity is not the class-giveaway.

Methodological rule \#3: \emph{ceiling-matching by noise injection is not sufficient}. Design tasks where the raw bag-of-words is uninformative so the encoder must discover higher-order structure.

\section{CC-15: Composite Objectives Interfere}

\textbf{Setup.} VQ + entropy anti-collapse, optionally with a 0.3-weighted $\Sn$-contrastive loss on the projection head.

\paragraph{Result.} VQ-only recovers purity 0.979. VQ+contrastive plateaus at purity 0.615.

\paragraph{Interpretation.} \emph{Destructive interference, not synergy.} The contrastive gradient pushes same-multiset-perms together in projection space, but it also pushes different-multisets apart ,  and the encoder's parameters are shared with the VQ commitment/codebook losses. The two objectives pull on the representation in incompatible directions, and VQ's ability to quantise into clean clusters is degraded.

This is a real finding for R's objective design: VQ and contrastive cannot be combined naïvely. If both are wanted, they likely need separate heads (a contrastive projection head and a VQ quantisation head on different transforms of the encoder output) or sequential training (pretrain with one, fine-tune with the other).

\section{CC-16: Codebook Over-Sizing Hurts Purity}

\textbf{Setup.} VQ-only with $K \in \{4, 8, 16, 32\}$ on the true-$K = 4$ task.

\paragraph{Result.} Purity and utilisation (Table~\ref{tab:cc16}).

\begin{table}[t]
\centering\small
\caption{CC-16. Purity vs codebook size on a $K=4$ task.}
\label{tab:cc16}
\begin{tabular}{rrr}
\toprule
Codebook $K$ & Purity & Clusters used \\
\midrule
4  & \textbf{0.979} & 4 \\
8  & 0.513 & 8 \\
16 & 0.734 & 14 \\
32 & 0.765 & 26 \\
\bottomrule
\end{tabular}
\end{table}

\paragraph{Interpretation.} With $K=8$ on a true-$K=4$ task, VQ splits each true class across two clusters, but the entropy regulariser pushes those splits toward equal utilisation ,  so purity drops as half of the per-cluster content is ``wrong'' (the partner half of the split class). At $K \in \{16, 32\}$, the entropy pressure weakens (uniform distribution over 16 or 32 slots is harder to achieve), and the model falls back toward a sparser representation, partially recovering purity.

For LORN's R: if the true latent class count is unknown, start with $K$ slightly smaller than the expected number, not larger. Over-sizing the codebook with entropy anti-collapse \emph{actively harms} structure discovery. This is a design constraint: R's codebook size is a hyperparameter that interacts strongly with the entropy weight.

\section{CC-17: First LORN Integration Test ,  Task-Design Failure}

\textbf{Setup.} VQ-R (4 clusters) + AdaLN bridge + attention-based L. Joint training on a conditional-prediction task: $\mathrm{target} = (\max\mathrm{token} + 3 \cdot \mathrm{class}) \bmod V$. Baseline: L alone with $\alpha = 0$.

\paragraph{Result.} LORN accuracy 0.943; L-alone accuracy 0.943; \textbf{lift over baseline = 0.000}. R's cluster purity 0.465 ,  R did not succeed at clustering under joint training.

\paragraph{Interpretation.} The task was solvable by L alone, so R contributed nothing. The intended conditional target was meant to require class information, but the class-signal is \emph{leaked} through the input: the maximum token of a class-$c$ sample is usually in the canonical block $[c \cdot 8, (c+1) \cdot 8)$, so $\max\text{tok}$ correlates strongly with class. L can infer class from the inputs alone and does not need R's $\alpha$.

R's cluster purity decaying to 0.465 is a secondary finding. Under joint training, R's gradient from the classification head is dominated by the downstream task's CE loss; the VQ commitment signal is a small side effect. Since L does not need R, R has no gradient incentive to cluster correctly ,  it just produces whatever $\alpha$ vector minimises the shared loss.

\paragraph{The methodology lesson.} A proper LORN integration test requires a task where L \emph{strictly cannot} succeed without R's signal. The canonical way to enforce this: restrict L's input view (e.g., L sees only 2--3 tokens at random per sample, not the full sequence) while R sees everything; then R must extract the class-signal from its full view and hand it to L through the bridge. Without this restriction, the bridge is redundant.

This is a design lesson, not a refutation. CC-17 has not shown that LORN integration fails; it has shown that integration testing requires a task where the bridge is \emph{information-essential}, not just information-helpful.


\section{Round 7 Summary and Outstanding Work}

\begin{itemize}[nosep]
\item \textbf{CC-0b}: Overlap task still probe-saturated. Need a task where raw bag-of-words is uninformative.
\item \textbf{CC-15 (kill, composite):} VQ + contrastive interfere destructively. Composite R objectives must use separate heads or sequential training.
\item \textbf{CC-16 (kill, over-sizing):} Larger codebook hurts purity under entropy regularisation. Start $K$ small.
\item \textbf{CC-17 (task-design failure):} L can solve the task without R. Next integration test requires an information-essential bridge.
\end{itemize}

\paragraph{Updated constraints after seven rounds.}
The seven architectural constraints from Round 6 stand. Round 7 adds three \emph{design} constraints:
\begin{enumerate}[nosep,leftmargin=*,resume]
\setcounter{enumi}{7}
\item \textbf{Composite R objectives need architectural separation.} VQ and contrastive cannot share a single encoder body without interference.
\item \textbf{VQ codebook size must match true $K$ or be smaller.} Over-sized codebooks degrade purity under entropy anti-collapse.
\item \textbf{Integration testing requires information-essential bridges.} L must be information-starved without R, or the bridge is mute.
\end{enumerate}

\paragraph{Round 8 targets.}
(a) \textbf{Information-essential task.} L sees only 2 tokens out of 12 per sample; R sees all 12. The class must be inferable from 12 but not from 2. Full LORN becomes the minimum sufficient architecture.
(b) \textbf{Probe-informative task.} Structured co-occurrence signatures (class $c$'s samples contain a specific token \emph{pair} more often; individual token frequencies are class-flat). Raw bag-of-words does not separate classes.
(c) \textbf{CC-17b}: re-run the integration test with the information-essential task.
(d) \textbf{CC-18}: ablate the bridge (AdaLN vs hypernet vs zero-$\alpha$) on the information-essential task, to measure bridge lift directly.


% ======================================================================
\part*{Round 8: The Integration Succeeds}

The first LORN end-to-end integration test with a task L genuinely cannot solve alone. Results are decisive: LORN lifts substantially over L-alone; one bridge design gets within 1\% of the oracle.

\section{The Information-Starved Task}

R receives 12 tokens sampled from class $c$'s Dirichlet distribution (55\% canonical block, 45\% uniform noise). L receives only 2 of those 12 tokens, chosen uniformly at random per sample. The target is
\[
t \;=\; (x_{L,0} + x_{L,1} + 7 \cdot c) \bmod V,
\]
so L needs both its two visible tokens \emph{and} the class $c$ to predict. With only 2 tokens, L cannot reliably infer $c$. With R's $\alpha$ carrying $c$, it can.

\paragraph{Baseline ceilings.}
\begin{itemize}[nosep]
\item \textbf{Oracle L} ,  L receives the \emph{true} class as a one-hot $\alpha$: 0.896. This is the achievable ceiling given 2-token visibility.
\item \textbf{L-alone} ,  L receives $\alpha = 0$: 0.642. This is what 2-token visibility alone can achieve. Substantially above chance $(1/V = 0.031)$ because $x_{L,0} + x_{L,1}$ itself carries information.
\end{itemize}

\section{CC-17b: LORN vs L-Alone vs Oracle}

\textbf{Setup.} LORN = VQ-R on 12-token input + AdaLN bridge + L on 2-token input. Jointly trained for 4000 steps. Three models trained in parallel: LORN, L-alone (zero bridge), Oracle-L (true class embedding as $\alpha$).

\paragraph{Result.} \emph{Pass. Substantial lift over L-alone, $\sim$77\% of the oracle gap closed.}

\begin{table}[t]
\centering\small
\caption{CC-17b. Full LORN lifts over L-alone on an information-starved task.}
\label{tab:cc17b}
\begin{tabular}{lr}
\toprule
Model & Val acc \\
\midrule
L-alone (zero-$\alpha$) & 0.642 \\
\textbf{LORN (VQ-R + AdaLN)} & \textbf{0.819} \\
Oracle (true class as $\alpha$) & 0.896 \\
\midrule
LIFT (LORN $-$ L-alone) & $+0.177$ \\
GAP (Oracle $-$ LORN) & $+0.077$ \\
R purity at step 4000 & 0.900 \\
\bottomrule
\end{tabular}
\end{table}

\paragraph{Training dynamics.} R's cluster purity climbs from 0.43 (random VQ init) to 0.90 over 4000 steps. L-alone plateaus at 0.642 by step 1000 ,  that is the 2-token-visibility ceiling. LORN continues improving as R's clusters sharpen, reaching 0.819. The full LORN model gets 77\% of the way from L-alone to the oracle, using an unsupervised cluster discovery as its only class signal.

\paragraph{Interpretation.} This is the integration proof that Round 7 was missing. The task is designed so that class-information is \emph{strictly necessary} for L; R discovers the class through VQ; the AdaLN bridge carries R's discovery to L; L uses it to close most of the oracle gap. Every architectural component has a measurable contribution.\footnote{\emph{Scope.} The +17.7-point lift shows that the bridge carries signal and that R's cluster assignment is usable by L under joint training. The R here uses mean-pool readout and the classes are bag-of-words-accessible. Per the Phase 1 checklist this does not count as evidence of structural capability; it is mechanism validation that the bridge works and that R does not collapse under joint training.}

The 0.077 residual gap to the oracle is the cost of using a 4-cluster VQ instead of the true 4-class labels ,  10\% of R's assignments remain wrong (purity 0.90), and those propagate to L's prediction error. Longer training or a more expressive R would likely close some of this.

\section{CC-18: Bridge Ablation}

\textbf{Setup.} Same R, same task, same training regime as CC-17b ,  only the bridge type varies: zero (no $\alpha$ signal), concat (add linear-projected $\alpha$ to hidden state), AdaLN (scale-shift modulation), hypernet (low-rank weight delta generated from $\alpha$).

\paragraph{Result.} Hypernet dominates. See Table~\ref{tab:cc18}.

\begin{table}[t]
\centering\small
\caption{CC-18. Bridge ablation on the information-starved task. $R_p$ = R's cluster purity under joint training.}
\label{tab:cc18}
\begin{tabular}{lrrr}
\toprule
Bridge & Val acc & $R_p$ & Params \\
\midrule
zero    & 0.637 & 0.982 & 608K \\
concat  & 0.685 & 0.518 & 641K \\
adaln   & 0.767 & 0.695 & 674K \\
\textbf{hyper}   & \textbf{0.887} & \textbf{0.950} & 905K \\
\midrule
Oracle  & 0.896 & -- & -- \\
\bottomrule
\end{tabular}
\end{table}

\paragraph{Interpretation.} Three striking facts:

\begin{enumerate}[nosep,leftmargin=*]
\item \textbf{Hypernet closes the oracle gap to 0.009.} The best bridge gives LORN within one percentage point of the oracle ceiling. The full architecture, trained end-to-end with \emph{no class labels}, is essentially optimal on this task.
\item \textbf{R's purity under joint training depends on the bridge.} Zero bridge (no pressure on R) gives clean VQ clusters (0.982 purity). Concat bridge \emph{collapses} R's clustering (0.518 ,  clusters contaminated) because concat's weak information channel forces R to distort its $\alpha$ to compensate. AdaLN is intermediate (0.695). Hypernet recovers \emph{near-zero}-pressure clustering (0.950) because it has enough bandwidth to accept R's cluster structure as-is without forcing distortion.
\item \textbf{Bridge type matters more than parameter count.} Hypernet has $\sim$32\% more parameters than AdaLN, but its advantage (+0.12) far exceeds what extra parameters elsewhere could buy.
\end{enumerate}

Architectural takeaway: hypernet is the right bridge for LORN \emph{at this toy scale}. It has enough capacity to transmit R's full cluster signal; it is low-pressure enough that R's discovery dynamics survive joint training; and it produces an end-to-end accuracy that essentially matches the oracle.\footnote{\emph{Scope.} This ablation compares bridge \emph{mechanisms} on the information-starved toy, with a mean-pool R and bag-of-words-accessible classes. It shows the hypernet channel carries more of R's signal than AdaLN or concat on this task. It does not establish that hypernet is the Phase-1 bridge of choice; that comparison has to be redone with a Phase-1-compliant R and a non-bag-of-words task.}


\section{Round 8 Summary ,  The Programme's Constraint Set is Settled}

\begin{itemize}[nosep]
\item \textbf{CC-17b (pass, +17.7 points):} LORN with AdaLN lifts 0.177 over L-alone, closing 77\% of the oracle gap. R's VQ clusters reach 0.90 purity under joint training.
\item \textbf{CC-18 (hypernet wins):} hypernet bridge reaches 0.887, within 0.009 of the oracle. R purity under hypernet is 0.950 ,  clustering survives joint training because the channel has enough bandwidth.
\end{itemize}

\paragraph{What this closes.}

Every architectural constraint in the LORN design specification now has end-to-end empirical backing at laptop scale:

\begin{enumerate}[nosep,leftmargin=*]
\item Composition = attention / bilinear (CC-3/4/4b kill; CC-9/11 positive).
\item R$\to$L bandwidth abundant at low $k$; \textbf{hypernet is the preferred bridge} (CC-5/6/17b/18).
\item Contraction trainable, $\lambda \approx 0.8$ (CC-7/8/8b).
\item Joint training eliminates chained burn-in (CC-10).
\item Attention discovers group structure automatically (CC-11).
\item Multi-horizon training stabilises variable-depth attractors (CC-8b).
\item Unsupervised R = attractor-basin codebook + anti-collapse (CC-13b).
\item Composite R objectives interfere ,  use separate heads or sequential (CC-15).
\item VQ codebook size must match or under-estimate true $K$ (CC-16).
\item \textbf{Integration works when L is information-starved and the bridge has enough bandwidth} (CC-17b/18).
\end{enumerate}

Ten constraints, twenty-three circuits, eight rounds. Every constraint has at least one pass-criterion positive result. Every kill corresponds to a specific architectural or task-design falsification. The battery has reached its natural stopping point at laptop scale: there is no circuit we can identify a priori that is not either resolved or out of laptop reach.

\paragraph{Scope caveat.}
The recipe that follows is the outcome of this laptop-scale battery. Every R tested in the battery uses a mean-pool readout, and every task is synthetic and bag-of-words-accessible; neither condition satisfies the Phase 1 checklist in \emph{lorn\_gauge} v3. The recipe below is therefore a \emph{starting point for the toy pipeline}, not the Phase 1 design, and it has not been validated at production scale. Treat ``VQ-R $\to$ hypernet $\to$ ORN-L'' as the stack that worked on these toys, pending a Phase-1-compliant re-evaluation.

\paragraph{The LORN recipe.}
The architecture is three pieces: R (observer), control (actor), L (predictor). These are distinct and must not be conflated.

\begin{quote}
\textbf{L-hemisphere (predictor):} any deep autoregressive model with a substantially order-independent residual mode. ORN is the L we tested; vanilla transformer, Mamba, Jamba, or any similar architecture should also work. ORN-specific results (Section 11 layer-27 recovery in \cite{gauge-v2}; CC-11 group discovery through shared-M) are \emph{supporting evidence}, not prerequisites.

\textbf{R-hemisphere (observer):} attention-based encoder + VQ codebook of size $K \lesssim$ expected latent class count + entropy anti-collapse regulariser. $\alpha$ = indexed codebook vector (straight-through). R is read-only: it watches L's inputs (or L's residual stream) and produces a structural fingerprint. R does not, by itself, act on L.

\textbf{Control (bridge):} the active component that uses R's $\alpha$ to change something about L. This is the piece whose purpose is still under active research (see below). Candidates include hypernet weight modulation (CC-18), AdaLN (CC-5/18), generation-time selection/filtering (Exp~F in \cite{gauge-v2}), steering-vector injection, cross-attention, or none of the above. The battery tested some options but did not settle which one is correct in a language-model setting.

\textbf{Training:} joint end-to-end with L's task-specific loss + small VQ commitment/codebook/entropy terms. For variable-depth inference components, sample training depth uniformly over $[h_{\min}, h_{\max}]$.
\end{quote}

\paragraph{Two distinct purposes for control, not yet disambiguated.}
The control piece can serve two fundamentally different purposes, and the circuit battery did \emph{not} settle which matters:

\begin{itemize}[nosep,leftmargin=*]
\item \textbf{Mode A ,  representation shaping (training-time co-evolution).} R's $\alpha$ is wired into L during training, so L's learned representation is better than it would be without R. CC-17b's +17.7-point lift lived in this mode. If L's structure is already good (as Exp~A in \cite{gauge-v2} suggests ,  74\% invariance floor emerges naturally), this may be unnecessary overhead.
\item \textbf{Mode B ,  behaviour steering (inference-time control).} L is trained independently. At inference, R's $\alpha$ influences L's generation: by selection/filtering over candidate outputs (Exp~F), by on-the-fly weight modulation, by steering vectors, or by prompt engineering. If L's structure is already good, this is the right path: R is a quality filter, not a teacher.
\end{itemize}

These are different control problems with potentially different winning architectures. The V4 strategy does not change ,  build and train at scale, measure ,  but the \emph{interpretation} of any result will hinge on which mode the control piece is operating in.

\paragraph{The real open problem: composing and generating from R to L.}
Beyond Modes A and B, there is a distinct unresolved question that neither the circuit battery nor the hemisphere sweeps have settled:

\begin{itemize}[nosep,leftmargin=*]
\item Given a chosen $\alpha^*$ in R-space, produce L-space text whose R-fingerprint matches $\alpha^*$. This is \emph{generation from structural intent}.
\item Given two $\alpha_1, \alpha_2$, compose them ($\alpha_1 + \alpha_2$, interpolation, or group operation) and generate text with the composed structure. This is \emph{R-space arithmetic manifesting in L-space}.
\item Navigate R-space continuously ,  move from explanatory to argumentative structure at fixed topic, from formal to casual register at fixed semantics, and so on ,  and have L's generation follow.
\end{itemize}

The circuit battery shows that R exists (CC-13b: purity 0.97), that attention discovers group structure on $\alpha$-manifolds (CC-11: 99.3\% equivariance), and that a control piece can lift L performance when information-starved (CC-17b/18). But $\alpha \to$ \emph{generated text} is a bigger problem than $\alpha \to$ classification accuracy, and only a V4 on real text will tell us how hard it actually is. Exp~F (best-of-$N$ generation with $\alpha$-based selection) is the cheapest candidate mechanism; hypernet modulation during generation is another; steering-vector injection is a third.

This recipe and its three open questions are what we carry forward into the next scale.

\paragraph{What happens next.}

The circuit-exploration phase of the LORN programme is now complete. Continuations:

\begin{itemize}[nosep,leftmargin=*]
\item \textbf{V4 prototype} at 5--10M params on real text (TinyStories, WikiText-2). First measurement of whether the recipe scales.
\item \textbf{Theory synthesis} ,  a single paper tying the skew--dissipative theorems, the monoid-of-layer-actions framing, and the ten empirical constraints together.
\item \textbf{Targeted follow-up circuits} ,  slow-feature R on a harder multiset task; hierarchical VQ (multi-scale codebooks); structured $\alpha$-manifolds (non-discrete R outputs); attention-based composition at larger $p$ to test the $\log_2 p$ Fourier scaling.
\end{itemize}

The circuit battery has done its work.


% ======================================================================
\part*{The Integrated Architecture (LORN v4)}

After eight rounds of circuits, the design specification is complete enough to instantiate as a single trainable module. This section documents the integration and the \emph{honest} scale-readiness assessment.

\section{What Is Built}

\texttt{orn/lorn\_v4.py} is a single Python file (~280 lines) wiring the three components validated by the battery:

\begin{itemize}[nosep,leftmargin=*]
\item \textbf{VQREncoder} (the R-hemisphere). Attention encoder (no positional embedding ,  invariance by mean-pool) $\to$ learned codebook of size $K$ $\to$ VQ-VAE straight-through estimator with commitment, codebook, and entropy-anti-collapse losses. Validated by CC-13b (purity 0.381 $\to$ 0.974) and CC-17b/CC-18 (survives joint training at purity 0.95 with hypernet bridge).
\item \textbf{HypernetBridge.} For each sample, generate a low-rank weight delta $\Delta W = U(\alpha) V(\alpha)^\top$ for the top-$n$ L-layers. Validated by CC-18: beats AdaLN by 12 points on the information-starved task and closes the oracle gap to 0.009.
\item \textbf{ORNlite} (the L-hemisphere reference implementation). Shared $\M = AB^\top$ coupling, shared wide FFN, per-layer gated contractive correction. Laptop-scale port of \texttt{orn/full\_orn.py}'s FullORN. For production scale, the HypernetBridge hooks into FullORN's blocks by modifying their forward to accept the per-sample $\Delta W$.
\end{itemize}

The joint training loop follows the recipe from the battery:
\begin{itemize}[nosep]
\item Task loss (CE on L's output) + $0.1 \cdot (\text{commit} + \text{codebook}) - 0.1 \cdot \text{entropy}$
\item If variable-depth inference is required anywhere, sample training depth uniformly over $[h_{\min}, h_{\max}]$ (CC-8b: eliminates over-iteration drift).
\item Codebook size $K \le$ expected latent class count (CC-16: over-sizing hurts purity).
\item Do not combine VQ with contrastive on a shared encoder body (CC-15 kill).
\end{itemize}

Smoke test at the bottom of \texttt{lorn\_v4.py} instantiates a 540K-parameter model, runs a forward pass and a joint backward pass, confirming end-to-end gradient flow.

\section{What Is Not Built}

\emph{Being honest about the gap between specification and deployment.}

\begin{itemize}[nosep,leftmargin=*]
\item \textbf{No production-scale trainer.} The \texttt{ORNlite} in \texttt{lorn\_v4.py} is a laptop reference. The production L lives in \texttt{orn/full\_orn.py} with RoPE, GQA, SDPA, AMP, weight-tying, and FP16/BF16 plumbing. We have \emph{not} yet wired \texttt{HypernetBridge} into \texttt{FullORN}'s blocks. Engineering task, not research risk: it needs a 20-line change in \texttt{ORNBlock.forward} to accept a per-sample delta.
\item \textbf{No real-text dataset integration.} The circuit tests use synthetic algorithmic tasks. LORN v4 has not been run on TinyStories, WikiText, or FineWeb. The data pipeline (\texttt{training/train\_v2.py}) is production-ready but does not know about R's $\alpha$ or the bridge. Another engineering task.
\item \textbf{No distributed / multi-GPU training.} The current tests run on MPS or a single GPU. Scale-up to RTX 4090 (the 605M target) requires only standard PyTorch DDP wrapping ,  no research risk.
\item \textbf{No inference API.} We have not defined the inference surface for LORN: does R re-encode at every generation step, or does $\alpha$ stay fixed for the generation? Does R see the prefix, or only the prompt? These are inference-path design choices we have not made.
\item \textbf{No downstream benchmarks.} ORN v3 has been benchmarked on 8-task lm-eval. LORN has no benchmark numbers at any scale because it has not been trained on real data yet.
\end{itemize}

\section{Scale-Readiness: Direct Answer}

\textbf{We have the specification and all the components. We have not yet built a single trainer that runs LORN at production scale on real text.}

The gap is entirely engineering. Concretely, to get from here to a trainable V4 at 5--10M parameters on TinyStories:

\begin{enumerate}[nosep,leftmargin=*]
\item Modify \texttt{orn/full\_orn.py} \texttt{ORNBlock.forward} to accept an optional per-sample delta tuple $(U, V)$ and apply it to $W_V$ or the SharedM coupling. (~20 lines.)
\item Add an \texttt{ORN\_with\_Bridge} class that wraps \texttt{FullORN} with an injected \texttt{HypernetBridge}. (~30 lines.)
\item Fork \texttt{training/train\_v2.py} to \texttt{training/train\_lorn.py}: accept a two-headed input (full sequence for R, partial or full for L), run R forward to get $\alpha$, pass deltas to L, compute joint loss. (~100 lines.)
\item Start with TinyStories; log R's cluster purity alongside val loss.
\end{enumerate}

Estimated wall time: one day of focused engineering for the integration. After that, a multi-day training run on a rented 4090 tests whether the recipe actually produces a useful language model.

\section{What We Expect}

Two honest hypotheses about the first scale-up run:

\begin{itemize}[nosep,leftmargin=*]
\item \textbf{Optimistic.} R discovers coarse topic / register clusters on natural text (like the 605M hemisphere sweep already showed qualitatively). The hypernet bridge modulates L's coupling in a way that reduces val loss relative to a matched-parameter ORN-alone. R + L are complementary, not redundant.
\item \textbf{Pessimistic.} The information-starvation that made CC-17b work (L sees only 2 of 12 tokens) does not map cleanly to a language-model trainer, where L sees the full prefix by construction. Without an information-essential bridge, R may get zero gradient (CC-17 failure mode). A proper scale-up may need a \emph{new} task design: for example, R's $\alpha$ modulates only a subset of L's layers, or only the L-output classifier head, or only the final layer's RMSNorm.
\end{itemize}

This ambiguity is itself a battery finding: \emph{the bridge is only worth its parameter budget if L needs R's information.} A language model with full autoregressive input may already have enough signal in the prefix. CC-17b's design ,  L sees 2 tokens, R sees 12 ,  is the small-scale analogue of something like: L is an autoregressive head, R is a long-context summariser, and L is forced to operate on a context-truncated view. A V4 that mirrors this architecture ,   long-context R + truncated-context L ,  would preserve the information-essentialness.


\section{Summary}

All ten constraints are empirically grounded. The unified module exists as \texttt{orn/lorn\_v4.py} and passes a forward/backward smoke test. The bridge from laptop-scale capability to production training is roughly one day of engineering plus a day of training time. The primary research risk at scale is whether the bridge remains information-essential in an autoregressive setting; the engineering risk is near zero.

The research programme is ready for its first scaled test.


% ======================================================================
\section*{Relation to LORN v3}

This paper was written before the \emph{lorn\_gauge} v3 Phase 1 checklist existed. Under that later framework, the scope of what this battery can and cannot claim is sharply delimited. We make that scope explicit here so the document can stay usefully on record as a mechanism archive without being mistaken for Phase 1 evidence.

\paragraph{What the battery \emph{does} show (mechanism validation).}
\begin{itemize}[nosep,leftmargin=*]
\item \textbf{The R$\to$L bridge carries signal at low bandwidth.} CC-5 (AdaLN saturates at $k{=}1$) and CC-6 (hypernet matches AdaLN) show that compact $\alpha$ channels have usable bandwidth for mode-selection tasks. CC-18 shows that hypernet transmits more of R's cluster structure than AdaLN or concat under joint training on this toy.
\item \textbf{Hypernet is a viable channel.} As a mechanism, a low-rank per-sample weight delta $\Delta W = U(\alpha) V(\alpha)^\top$ closes the oracle gap on the information-starved toy to 0.009 and leaves R's clustering near-intact (purity 0.95).
\item \textbf{Contractive residual blocks exhibit measurable $\lambda$.} CC-7 ($\lambda = 0.844$) and CC-8 ($\lambda = 0.777$) show geometric decay to the attractor set; CC-8b shows multi-horizon training eliminates over-iteration drift. These are clean, quantitative measurements of contraction as an architectural property.
\item \textbf{Attention on symbolic groups is equivariant.} CC-9 (97.7\% on Z$_{17}$ addition) and CC-11 (99.3\% cyclic-shift equivariance) show attention-based composers discover the group structure of small abelian groups without explicit equivariance training. This is evidence about attention-on-groups, not about structural capability on language.
\item \textbf{Joint training with VQ-R does not collapse R on these synthetic tasks.} CC-17b/CC-18 show R's cluster purity survives joint training with an information-essential bridge, provided the channel is wide enough (hypernet $>$ AdaLN $>$ concat).
\end{itemize}

\paragraph{What the battery \emph{does not} show.}
\begin{itemize}[nosep,leftmargin=*]
\item \textbf{That R captures structure beyond bag-of-words.} Every task here is synthetic with class-distinguishing tokens accessible from the raw multiset. A mean-pool encoder is sufficient on construction. Phase 1 requirement \#14 explicitly retires such labels as evidence of structural capability.
\item \textbf{That the architectures scale.} All tests run at $\lesssim 1$M parameters on Apple M4 in under 20 minutes. Scaling behaviour is unmeasured.
\item \textbf{That composition works for non-symbolic content.} CC-9/CC-10/CC-11 concern Z$_{17}$ and $S_4$. Attention discovers the group structure of these closed algebraic systems; no claim follows for linguistic composition, analogical composition, or any non-symbolic task.
\item \textbf{That any of this reaches the Phase 1 checklist's structural-evaluation bar.} The R architectures tested here all use mean-pool readouts (disqualified by checklist requirement \#4). The tasks are all bag-of-words-accessible (disqualified by requirement \#14). No result in this document speaks to the Phase 1 question of whether LORN's R captures non-bag-of-words structure at scale.
\end{itemize}

\paragraph{How to read the rest of the paper.}
Every verdict above ,  pass, kill, partial ,  is a statement about the mechanism at laptop scale on toy data. The ``preferred stack'' recommendation and the ``LORN v4'' recipe are pipeline fixtures useful for running further mechanism tests, not architectural commitments. The Phase 1 question ,  \emph{does R capture structure a bag-of-words baseline cannot?} ,  is not addressed anywhere in this document. That question is the subject of \emph{lorn\_gauge} v3 and its successor experiments.


% ======================================================================
\begin{thebibliography}{99}
\bibitem{gauge-v2} Sirsh (2026). Two Hemispheres for One Language Model: Permutation-Invariant Structure via the Orbital Response Network.
\bibitem{building-circuits} Sirsh (2026). Building Circuits: A Tutorial on Representational Capacity Tests, and a Programme for the Gauge LORN.
\bibitem{zaheer2017deepsets} Zaheer, M., Kottur, S., Ravanbakhsh, S., P\'{o}czos, B., Salakhutdinov, R., Smola, A. (2017). Deep Sets. NeurIPS.
\bibitem{peebles2023dit} Peebles, W., Xie, S. (2023). Scalable Diffusion Models with Transformers. ICCV.
\bibitem{oord2018infonce} van den Oord, A., Li, Y., Vinyals, O. (2018). Representation Learning with Contrastive Predictive Coding.
\bibitem{chen2018neuralode} Chen, R. T. Q., Rubanova, Y., Bettencourt, J., Duvenaud, D. (2018). Neural Ordinary Differential Equations. NeurIPS.
\end{thebibliography}

\end{document}
